\subsection{COIE Schemes }
We assume the input index vector $v \in \zo^n$ is sparse. In particular,
throughout the paper, we assume $s = o(n)$.

\subsubsection{A Warm-up construction}
Using an algebraic BF, we can create an $(n, s, c, f_p)$-COIE scheme
(the parameters $c$ and $f_p$ will be worked out after the description of
the scheme).


\paragraph{$\Encode(\lbra v_1 \lket, \ldots, \lbra v_n \lket)$.}
The encoding algorithm works as follows:

\BEN
\item Initialize a BF $\lbra B \lket := (\lbra B_1 \lket, \ldots, \lbra
B_c \lket)$ with $B_j = 0$ for all $j$. Let $\HHH = \{ h_q: \zo^* \to
[c] \}_{q=1}^\eta$ be the associated hash functions. 

\item For $i = 1, \ldots, n$:
    \BEN
    \item For $q = 1, \ldots, \eta$, do the following: 
      Compute $j = h_q(i)$ and set $\lbra B_j \lket := \lbra B_j \lket +
      \lbra v_i \lket$.
 
\EEN
\EEN



Note that at step 2.a in the above, if $v_i = 0$, then $B_j$ stays the
same. On the other hand, if $v_i = 1$, then $B_j$ will be increased by
1.  This implies that $B$ will exactly store the results of the
operations $\{ {\sf BF.Insert}(B, i) : i \in \idx(v)\}.$


\paragraph{$\Decode(B_1, \ldots, B_c)$.}
Given the algebraic BF $B$, we can recover
the indices for the nonzero elements as follows:
\BI
\item Initialize $I$ to be the empty set. 
\item For $i \in [n]$: if ${\sf BF.Check}(B, i)$ = ``yes", add $i$ to $I$.  
\item return $I$. 
\EI

\paragraph{Parameters $c$ and $f_p$.} Since this is a warm-up
construction, we perform only a rough estimation on the false positive parameter
and the compactness parameter. 

For reasons that will become clear later, we wish to keep the upper bound
on the number of false positives ($f_p$) small.  In
particular, we use a BF with false-positive rate $1/n$. Since there
are $n$ operations of $\sf BF.Check$, the expected number of false
positives is 1, and from the Chernoff bound, the number of false
positives is bounded by $\Omega(\log \secparam)$ with overwhelming
probability in $\secparam$. This implies that we have $f_p = \Omega(\log
\secparam).$

The dimension $c$ of the Bloom filter $B$ can be
computed using the following equation of BF false positive ratio: 

$$ \Big(1 - e^{-\frac{\eta s}{c}} \Big)^\eta \le \frac 1 n,$$ 

Setting $c = \eta s \cdot n^{\frac 1 \eta}$ will satisfy the
equation. This can be verified by using an equality $1 - e^{-x} \le x$
for $x \in [0,1]$; that is, $1 - e^{-\frac{\eta s}{c}} \le \frac{\eta
s}{c} = 1/n^{1/\eta}.$

\paragraph{Efficiency.}
\BI
\item The encoding algorithm uses $n\eta$ homomorphic addition
operations, and $n\eta$ hash functions.  

\item The dimension $c$ of the encoding is $\eta s \cdot n^{\frac 1
\eta}$. Usually, $\eta$ is set to between 2 and 32.

\item The decoding algorithm uses $n$ operations of $\sf BF.Check$.
\EI

In summary, we have reduced the encoding size $c$ to be sub-linear in $n$ as desired. However, we still need to reduce the number $\sf
BF.Check$ operations in Decode to be sub-linear in $n$. We show how to achieve that in our next construction.

\subsubsection{BF-COIE}\label{sec:BF-COIE}
We now show how to improve the above construction to achieve decoding in
time $o(n)$.  The main idea of this improvement is to use Bloom filters
to represent a binary search tree, one BF per level of the tree.  We can
then guide the decoding algorithm to avoid decoding branches that do not contain
non-zero entries.  As most branches can be truncated well before
reaching the leaf-level Bloom filter, this results in sublinear total
cost.

% By improving the above warm-up scheme, we will construct a BF-cased COIE
% scheme whose decoding algorithm performs $o(n)$ operations of $\sf
% BF.Check$. The main idea is to {\em guide the decoding algorithm to
% perform the binary search in parallel only for the necessary portion of
% indices}.


\paragraph{Example.} 
Before presenting the formal protocol for this construction we convey our idea through an example.  Let $n =
32$, and suppose we wish to encode the indices $I =  \{1, 15, 16\}$. 
Denote $$I^k = \left \{ \Big \lceil \frac i {2^k} \Big \rceil : i \in
I \right \}.$$

Intuitively, an element $i$ in $I^k$ can be thought of a range of length
$2^k$ covering $[(i-1)\cdot 2^{k}+1, i\cdot2^k]$. We have:
\BI
\item $I^4 = \{1\}.$
\item $I^3 = \{1, 2\}.$
\item $I^2 = \{1, 4\}.$
\item $I^1 = \{1, 8\}.$
\item $I^0 = \{1, 15, 16\}.$
\EI

Now, assume we insert each set $I^k$ into its own BF. We can traverse these BF's to decode the set $I$ as follows:

\BEN
\item Check $I^4$ for all possible indices. The only possible indices at this level are $1$ and $2$, since $n=32$ and $I^4$ divides the original indices by
$2^4 = 16$. 

In the above example, When we query the BF for $I^4$, it 
only contains the index $1$, which means that no values greater than 16 are contained in $I$.  We can thus avoid checking any such indices at the lower levels.  

Now consider the BF at the next level (i.e., the BF for $I^3$).  The
only possible values at this level are 1,2,3,4, but since we already
know that there are no values greater than 16 in $I$, we only need to
check for values $1, 2$ (since $3 \cdot 8 > 16$).

\item Check $I^3$ for indices $1, 2$. The BF will show that indices $1$ and $2$
are both present, which means that we need to check indices $1,2$ and $3,4$ in $I^2$. 
\item Check $I^2$ for indices $1, 2, 3, 4$. The BF will show that indices $1$ and $4$
are present, which means that we only need to check indices $1, 2$ and $7, 8$ in $I^1$, all other indices can be skipped.

\item Check $I^1$ for indices $1, 2, 7, 8$. The BF will show that indices $1$ and $8$
are present, which means that we need to check indices $1, 2$ and $15, 16$.
\item Check $I^0$ for indices $1, 2, 15, 16$, and output the
final present indices $1, 15, 16$. 
\EEN 

Assuming, for now, that there are no false positives, observe that this approach checks at most $2 \cdot |I|$ values at each
level, and there are $\lg n$ levels. Therefore, the decoding algorithm
will check $O( |I| \cdot \lg n)$ indices, which is sub-linear in $n$.  

\paragraph{BF-COIE.} We now describe our BF-COIE construction. As
before, we will work out the parameters after describing our
construction. The encoding algorithm is described in Algorithm~\ref{alg:BFCOIEE}.   

\begin{algorithm}[htbp]

\BEN
\item[] \hspace{-2em} For simplicity, $n$ and $s$ are assumed to be powers of 2. 

\item $t := \lg {\frac n {2s}}$ 
\item For $k = 0, \ldots, t$: 

\BEN
\item Initialize $\lbra B^k \lket = (\lbra B^k_1 \lket, \ldots, \lbra
B^k_\ell \lket) := (\nil, \ldots, \nil)$. 

\item Choose $\HHH^k = \{ h^k_q: \zo^* \to [\ell] \}_{q=1}^\eta$ at random.

\item For $i \in [n]$ and for $q \in [\eta]$: \\ 
  \BEN
  \item[] $i' := \lceil i/2^k \rceil$, 
        $j := h^k_q(i')$, \\
        If $\lbra B^k_j \lket$ is $\nil$, then $\lbra B^k_j \lket := \lbra v_{i'} \lket$ \\
        Otherwise, $\lbra B^k_j \lket := \lbra B^k_j \lket + \lbra v_{i'} \lket$
  \EEN
\EEN

\item Output $\lbra B^0 \lket, \ldots, \lbra B^t \lket$.
\EEN

\caption{ ${\sf BF\mbox{-}COIE}.\Encode(\lbra v_1 \lket, \ldots, \lbra v_n \lket)$}
\label{alg:BFCOIEE}
\end{algorithm}

Note that in steps (a) to (c) above, the warm-up construction is used to
construct BF $B^k$ for indices $I^k$. 

In order to reduce the size of the output encoding, we set $t$ to be
$\lg {\frac n {2s}}$ instead of $\lg n$ as described previously. Note
that when $t$ is set in this way, $I^t$ contains at most $n/2^t = 2s$
possible values thus maintaining our invariant.


% Note that with $B^t$, index
% $i'$ can still be at most $i/2^k = 3s$, which is an important aspect to
% set  the size $\ell$ of each BF below, when we work out the parameters. 

The decoding algorithm is described in Algorithm~\ref{alg:BFCOIED}.   

\begin{algorithm}[htbp]
\BEN
\item Initialize $I, I^0, \ldots, I^{t-1} := \emptyset$
\item Initialize $I^t := \{1, \ldots, n/2^t\} = [2s]$
\item For $k = t, t-1, \ldots, 1$, and for $i' \in I^k$: \\
  $~~~$ If ${\sf BF.Check}(B^k, i')$ is ``yes", add $2i'-1$, $2i'$ in $I^{k-1}$  
\item For $i \in I^0$: \\
  $~~~$ If ${\sf BF.Check}(B^0, i)$ is ``yes", add $i$ to $I$  
\item Output $I$ 
\EEN
\caption{ ${\sf BF\mbox{-}COIE}.\Decode(B^0, \ldots, B^t)$} 
\label{alg:BFCOIED}
\end{algorithm}

\paragraph{Useful lemma.}
The following lemma will be useful to analyze the parameters $c$ and
$f_p$.

\begin{lemma}\label{lem:4_1}
Consider a Bloom filter with false positive rate $\frac{1}{m}$, where
$m$ is an arbitrary positive integer. Suppose at most $m$ $\sf BF.Check$ operations
are performed in the BF.
%
Then, for any $\delta > 0$, we have:
$$\Pr [ \mbox{\rm \# false positives} \ge 1+\delta] 
  \le {\frac {e^\delta} {(1+\delta)^{(1+\delta)}}}.$$ 
\end{lemma}

The proof, by an application of the Chernoff bound, can be found in Appendix~\ref{app:proof41}. 
\iffalse
\begin{proof}
Let $\alpha_i$ be the $i$th item that is checked through $\sf
BF.Check$. That is, we consider a sequence of $${\sf
BF.Check}(\alpha_1), \ldots, {\sf BF.Check}(\alpha_{m}),$$
where $\alpha_i$ is an arbitrary item. 
Since we wish to upper bound the false positives (i.e., we don't care about true
positives), it suffices to consider the
case that for every $i$, 
$\alpha_i \not \in {\sf BF}$ (i.e, $\alpha_i$ has not been inserted in the
BF) as this maximizes the number of possible false positives.

Let $X_1, \ldots, X_m$ be independent Bernoulli random variables with
$\Pr[X_i = 1] = 1/m$. 
%
Since the BF false positive rate is assumed to be $1/m$, we have
for all $i$, 
$$\Pr[{\sf BF.Check}(\alpha_i)=1] = \Pr[\mbox{query $i$ is a false positive}] \le
1/m.$$
Thus, we can bound the number of false positives
by $\sum_{i=1}^m X_i$.

Now, let $\mu :=
\Exp[\sum X_i] = m \cdot \frac 1 m = 1$.  By applying the Chernoff bound with $\mu =
1$,  we have: 
%
$$ \Pr\left[\sum_{i=1}^m X_i \ge 1 + \delta \right]  
\le {\frac {e^\delta} {(1+\delta)^{(1+\delta)}}} .$$
\end{proof}
\fi

Regarding the above Lemma, we remark that setting $\delta =
\Omega(\log \secparam)$, we have $$\Pr\left[\sum_{i=1}^m X_i \ge 1 +
\delta \right]  = \negl(\secparam).$$


\paragraph{Parameters $c$ and $f_p$.} We set the false positive
upperbound $f_p := \Omega(\log \secparam)$ for the BF-COIE scheme. In
our experiments, we set $f_p = 16$. 

Now, let $m = \max(2s, s+2f_p)$, we set the BF false positive rate to $1/m$.
Recall that in the BF-COIE construction, the topmost BF $B^t$ performs the $\sf
BF.Check$ operation with $2s$ times; see line (2) in
Algorithm~\ref{alg:BFCOIED}. Using the above Lemma, the number of false
positives in the top level BF $B^t$ is at most $f_p$  with all but
negligible probability in $\secparam$. Furthermore, the index $i$ in
$B^t$ is expanded into two indices $2i-1$ and $2i$ in $B^{t-1}$. This
means that the number of false indices to be checked in $B^{t-1}$ due to
the false positives in $B^t$ is at most
$2f_p$. 

Now consider an index $i$ that belongs to $B^{t}$.  Algorithm~\ref{alg:BFCOIED} will 
run ${\sf BF.Check}$ on the values $2i-1$ and $2i$ in $B^{t-1}$.  Since at least one
of these values must actually belong to $B^{t-1}$, this leads to at most one false index
being checked.
%
Thus, the maximum number of false indices that would be checked in
$B^{t-1}$ is at most $s + 2f_p$ (i.e., $2f_p$ from false
positives of $B^t$ and $s$ from true positives of $B^t$). 

The above argument applies inductively all the way to the bottom most
level, which means that the maximum number of false indices that would
be checked in each level BF $B^i$ will be at most $s + 2f_p$.  In the
end, the bottom BF will have at most $f_p$ false positives, and the
overall BF-COIE scheme will have at most $f_p$ false positives with all
but negligible probability in $\secparam$. 
 
For the compactness parameter $c$, we must determine the dimension $\ell$
of each BF. Recall that we set the BF false positive rate to 
$1/m$ for $m = \max(2s, s+2f_p)$:

$$ \Big(1 - e^{-\frac{\eta s}{\ell}} \Big)^\eta \le \frac 1 m.$$ 
Setting 
$\ell = \eta \cdot s \cdot m^{\frac 1 \eta}$ would
satisfy the above condition, which can be  
verified using an inequality $1 - e^{-x} \le x$ for $x \in [0,1]$; that
is, $1 - e^{-\frac{\eta s}{\ell}} \le \frac{\eta s}{\ell} =
(1/m)^{1/\eta}.$

Since the encoding has $t+1$ BFs, the overall compactness parameter is as
follows:
$$c = (t + 1) \cdot \ell = O\left ( \eta \cdot s^{1 + \frac 1 \eta}
\cdot \lg \frac n s \right ).$$

\paragraph{Efficiency.}
\BI
\item The size $c$ of encoding is $O\left ( \eta \cdot s^{1 + \frac 1 \eta} \cdot
\lg \frac n s \right )$. In our experiment, we choose $\eta = 2$. 

\item The encoding algorithm uses $O(\eta
\cdot n \cdot \lg \frac n s)$ homomorphic addition operations and hash
functions. 

\item The decoding algorithm uses $\sf BF.Check$ operations for $O(s \lg
\frac n s)$ times. 
\EI

In summary, assuming $s = o(n)$, we reduced the encoding size $c$ to be
sub-linear in $n$.  Moreover, we also reduced the number $\sf BF.Check$
operations to be sub-linear in $n$. 

\paragraph{Remark.} Although this scheme has multiple BFs, the size of
encoding $c$ is smaller than that of the warm-up scheme! This is because
with multiple levels of BFs, we can relax the false positive ratio for
each BF. The encoding computation time was increased by a multiplicative
factor of $\lg \frac n s$. 

